A particle locked in a small region cannot be at rest: its momentum is at least about $p \approx \Delta p = \frac{h}{4\pi\,\Delta x}$.

\begin{abcliste}
\abc An electron is confined to a region of $\Delta x = \xeO$. Estimate its smallest momentum and its smallest kinetic energy in \si{\electronvolt}.
\abc A proton is confined to a nucleus, a region of $\Delta x = \xpO$. Estimate its smallest momentum and its smallest kinetic energy in \si{\electronvolt}.
\abc The region is made half as large. How does the smallest kinetic energy change?
\end{abcliste}
